% displayname : Cinématique des fluides
\documentclass[12pt,fancy]{cours}
\Chapit{4}{1}
\begin{document}
\maketitre{
Nous abordons dans ce chapitre les notions élémentaires utiles à la description de l'écoulement d'un fluide (liquide ou gaz). Quelques écoulements particuliers importants sont abordés. On ne s'intéresse pour l'instant pas aux causes qui régissent l'écoulement (voir chapitres suivants sur la dynamique).
}
{
\item Cours Statique des fluides - 0
% \item Documents de cours Statique des fluides - 2
\item TD Statique des fluides
}


%%%%%%%%%%%%%%%%%%%%%%%%%%%%%%%%%%%%%%%%%%%%%%%%%%%%%%%%%%%%%
\section{Modélisation d’un fluide}

\begin{definition}[Fluide]\label{def:}
Un \imp{fluide} est un état de la matière dans lequel un corps peut s'écouler et prendre la forme de son contenant.
\end{definition}
% On appelle \imp{fluide} tout état de la matière dans lequel un corps peut s'écouler, changer de forme et épouser celle de son contenant.
 On distingue qualitativement deux types de fluide suivant l'ordre de grandeur de la masse volumique.
 \begin{liste}
  \item Les \imp{liquides} sont des fluides peu compressibles, de masse volumique élevée ($\rho \sim \SI{e3}{kg.m^{-3}}$ pour l'eau).
  \item Les \imp{gaz} sont des fluides compressibles, de masse volumique faible ($\rho \sim \SI{1}{kg.m^{-3}}$ pour l'air dans les conditions ambiantes).
 \end{liste}
 Cette distinction peut être arbitraire, car on peut passer continûment d'un gaz à un liquide.
%  les gaz (comme l'air), de faible masse volumique $\rho \sim \SI{1}{kg.m^{-3}}$ dans les conditions ambiantes, et fortement compressible, des liquides (comme l'eau liquide), de grande masse volumique et peu compressible. 
 L'étude de l'écoulement des fluides est l'objet de la mécanique des fluides.
\subsection{Hypothèse des milieux continus}


Un fluide peut se décrire suivant trois échelles caractéristiques.
\begin{liste}
  \item 
  % \element[.]
  {Échelle macroscopique}
  L'échelle macroscopique $L$ correspond à la taille de l'écoulement (diamètre d'un filet d'eau, d'un obstacle). À cette échelle, on décrit le milieu matériel par un petit nombre de propriétés intensives comme en thermodynamique (masse volumique moyenne $\rho$, vitesse moyenne $\Vec{v}$, etc.). Ces propriétés globales ne permettent pas de connaître les détails de l'écoulement.
  % \item 
  % \element[.]
  {Échelle microscopique}
  L'échelle microscopique $\ell$ est déterminée par la distance typique entre atomes ou molécules. À cette échelle, les notions thermodynamique de température, pression, etc. n'ont pas de sens. De plus, il y aurait un nombre faramineux de particules à étudier, de l'ordre de $\Na \simeq 10^{23}$.
  %, on encore par le libre parcours moyen.
\begin{figure}[H]
  \centering
  \includegraphics[scale=1.1]{model_cont.pdf}
  \caption{Dans l'hypothèse des milieux continus, il existe une échelle mésoscopique $a$, intermédiaire entre les échelles microscopique $\ell$ et macroscopique $L$. Une particule de fluide (carré centré en M) est un volume de taille mésoscopique $a$. Les propriétés physiques des particules de fluide définissent les propriétés locales de l'écoulement, comme la vitesse  $\Vec{v}(M,t)$.}
  \label{fig:model_cont}
\end{figure}
%%%%%%%%%%%%%

% \element[.]
{Échelle mésoscopique}
Pour décrire les propriétés locales de l'écoulement, on définit la notion de \imp{particule de fluide} : il s'agit d'un élément de volume mésoscopique du fluide, de taille typique $a$ suffisamment petite devant $L$ pour qu'on puisse l'assimiler à un point matériel $M$, et suffisamment grande devant $\ell$ pour contenir un grand nombre de particules, de sorte à pouvoir définir les grandeurs intensives masse volumique locale $\rho(M)$, vitesse locale $\Vec{v}(M)$, etc. La particule de fluide est un modèle qui permet de décrire simplement les propriétés de l'écoulement. Elle ne correspond pas à une particule réelle de l'écoulement !
% Une particule de fluide contient ainsi un nombre suffisamment grand de particules pour pouvoir définir ses propriétés intensives, mais suffisamment petit pour assimiler la particule de fluide à un point matériel.
Si cette échelle mésoscopique $a$ existe, le fluide tout entier peut alors se décrire par un ensemble continu de particules de fluide : on dit que le fluide est un \imp{milieu continu}.
\end{liste}


\begin{prop}[Hypothèse des milieux continus]\label{prop:milieu_cont}
Il existe une échelle \imp{mésoscopique} $a$ intermédiaire entre les échelles microscopique $\ell$ (atomes, molécules) et macroscopique~$L$ (écoulement, obstacle), soit 
\begin{equation*}
\ell \ll a \ll L.
\end{equation*}
On peut alors décrire le fluide par un ensemble continu de \imp{particules de fluide} de taille $a$.
\end{prop}

% \noindent Cette échelle intermédiaire permet de conserver des informations locales sur les propriétés du fluide sans avoir à considérer les détails microscopiques.

\begin{exemple}{Ordres de grandeurs}
  L'hypothèse des milieux continus est-elle valable dans les cas suivants ?
\begin{quest}
  \item Filet d'eau liquide en sortie d'un robinet, \Fig{model_cont}. Dans les conditions ambiantes, la  densité volumique en molécules de l'eau liquide vaut $n \simeq \SI{3e28}{m^{-3}}$.
  \item Navette spatiale à la limite de la stratosphère, à $\SI{51}{km}$ d'altitude, où l'air se comporte comme un gaz parfait à la pression $p = \SI{66}{Pa}$ et la température $T = \SI{250}{K}$.
\end{quest}
\end{exemple}
\examplesol{
\setcounter{quest}{0}
\begin{quest}
  \item
  Le filet d'eau est de diamètre $L =\SI{1}{cm}$.
  La distance moyenne entre molécules d'eau est donc $\ell = n^{-1/3} \simeq \SI{0.3}{nm}$. Cette échelle est très petite comparée à la largeur du filet d'eau $L$. Il est donc possible d'y définir une échelle mésoscopique telle que $\ell \ll a \ll L$, par exemple $a \sim \SI{1}{\mu m}$.
  \item
   La distance entre particules de gaz de l'ordre de $\ell \simeq \SI{1}{m}$. L'écoulement de l'air autour d'une navette spatiale, de diamètre typique $L = \SI{2}{m}$ n'est donc pas bien décrit par un milieu continu : on ne peut pas y définir une échelle mésoscopique $a$ et une particule de fluide. Dans ce cas, on traite le milieu comme un système de points matériels. Ce régime pour lequel $\ell \sim L$ se nomme régime de Knudsen.
\end{quest}
}




\subsection{Champ des vitesses}

\parag{Description lagrangienne et eulérienne} Pour aborder l'étude des écoulements, on peut adopter deux points de vue.
\begin{liste}
  \item \udl{Lagrangien} : on suit une particule de fluide $\Pf$ et on décrit ses propriétés en fonction du temps $t$, comme on le ferait en mécanique du point. 
  La particule de fluide est alors caractérisée par sa température $T_{\Pf}(t)$, sa pression $p_{\Pf}(t)$, sa masse volumique $\rho_{\Pf}(t)$, sa vitesse $\Vec{v}_{\Pf}(t)$, etc, en fonction du temps $t$.
  \item \udl{Eulérien} : on suit l'évolution des propriétés physiques en un point $M$ de l'espace dans le temps $t$. 
  
  Un écoulement se caractérise ainsi par des champs scalaires : température $T(M,t)$, pression $p(M,t)$, masse volumique $\rho(M,t)$ ; ou vectoriels : vitesse de l'écoulement $\Vec{v}(M,t)$. Ce n'est donc pas la même particule étudiée aux différents points de l'espace et aux différents instants.
\end{liste}
  
Il est plus naturel dans un fluide d'adopter une approche \imp{eulérienne}.
C'est par exemple l'approche utilisée par les météorologistes pour décrire les champs de température, pression et vitesse du vent (Fig.~\ref{fig:carte_T}).  Dans le cadre du programme, nous adopterons uniquement cette approche.


\begin{figure}[H]
  \centering
  \subfigure[$t=\,$après-midi]{\label{fig:carte_T_1}
  % \includegraphics[height=0.4\textwidth]{champ_T_moy.jpg}
  \includegraphics[height=0.29\linewidth]{meteo_T_1.png}
  }\subfigure[$t=\,$soirée]{\label{fig:carte_T_2}
  \includegraphics[height=0.29\linewidth]{meteo_T_2.png}
  }\subfigure[$t=\,$nuit]{\label{fig:carte_T_3}
  \includegraphics[height=0.29\linewidth]{meteo_T_3.png}
  }
  \subfigure[]{\label{fig:carte_P}
  \includegraphics[height=0.4\textwidth]{vents.png}
   }
  %
  %\sousfigure{champ_T_moy.jpg}{fig:carte_T}{0.47}{vent_bretagne.pdf}{fig:carte_P}{0.5}
  \caption{(a-c) Température locale $T(M,t)$ pour plusieurs points $M$ en France à trois instants $t$ : (a) l'après-midi, (b) en soirée, (c) la nuit. Source : Météociel.fr. Un champ peut varier spatialement et temporellement. (d) Exemple de champ vectoriel $\Vec{v}(M,t)$ de vitesse du vent en France. Source : Windfinder.}
  \label{fig:carte_T}
  \end{figure}

  
% \begin{exemple}{Cartes météorologiques}
% Supposons que lors d'un trajet entre Tarbes et Paris, je relève à tout instant l'heure qu'il est $t$, ma position $M(t)$, et enfin la température $T$. J'obtiens une courbe de température $T(t)$ dont la variation a deux causes: d'une part mon déplacement entre Paris et Tarbes, sachant qu'à un instant donné il ne fait pas la même température dans ces deux villes ; d'autre part l'avancement de l'heure de la journée, sachant que la température augmente le matin puis diminue le soir. Une telle mesure illustre d'un point de vue lagrangien. Les météorologistes privilégient le point de vue eulérien : je place des thermomètres en de nombreux point $M$ en France, et je suis l'évolution de la température locale $T(M,t)$ dans le temps (Fig.~\ref{fig:carte_T}).

% \end{exemple}

\parag{Champ de vitesse}
Le vecteur $\Vec{v}(M,t)$ qui décrit la vitesse d'une particule de fluide située en $M$ à l'instant $t$ est appelé champ des vitesses. La représentation du champ de vitesse dans l'espace est appelée \imp{carte de champ}. La carte de champ s'apparente à une photographie (une représentation à un instant donné) de l'ensemble de l'écoulement (Fig.~\ref{fig:carte_P}). 
%Autrement dit ce sont les  eulérien $\Vec{v}(M,t)$. Elles sont définies mathématiquement par

\begin{definition}[Champ des vitesses]\label{def:}
  Le \imp{champ des vitesses} est le champ eulérien $\Vec{v}(M,t)$ qui représente la vitesse d'une particule de fluide au point $M$ à l'instant $t$. 
\end{definition}
\begin{remarque}
  Le champ des vitesses dépend du référentiel $\Rcal$ ! Nous nous placerons souvent dans le référentiel où les obstacles sont fixes (laboratoire, rive d'un fleuve, etc.).
\end{remarque}
Les lignes de champ du champ des vitesses sont appelées \imp{lignes de courant}. Un ensemble de lignes de courant qui s'appuient sur une section d'entrée et une section de sortie. est appelé \imp{tube de courant}.
\begin{definition}[Lignes et tube de courant]\label{def:}
  Les \imp{lignes de courant} sont les lignes de champ du champ de vitesse, i.e. les courbes en tout point tangentes et de même sens que le vecteur vitesse $\Vec{v}(M,t)$. 


  Un \imp{tube de courant} est un ensemble de lignes de courant qui s'appuient sur une section d'entrée et une section de sortie.
  
%   Elles vérifient
% \begin{equation*}
% \Vec{v}\pv \Vec{\d l}=\Vec{0},
% \end{equation*}
% avec $\Vec{\d l}$ le vecteur déplacement élémentaire.
\end{definition}
Dans le cadre du programme, la détermination analytique des lignes de courant n'est pas à connaître. Cependant, lorsque le champ des vitesses est porté par un vecteur de base, il est clair que les lignes de courant sont tangentes à ce vecteur de base. Il faut alors être capable de représenter les lignes de courant dans un tel cas. Dans les autres cas, on pourra s'appuyer sur un tracé numérique.


\begin{remarque}
  Il existe diverses méthodes expérimentales de visualisation des écoulements et de mesure de vitesse. Pour la visualisation des écoulements : injection de colorant, de bulles, ou de fumée dans l'écoulement (Fig.~\ref{fig:flow}), visualisation des variations d'indice par la méthode Schlieren, visualisation à partir de particules réfléchissantes ; pour la mesure de vitesse : anémométrie à fil chaud, vélocimétrie laser Doppler, PIV (Particle Image Velocimetry).
\end{remarque}

% \begin{methode}[Détermination des lignes de courant]
% \item Compter le nombre de composantes du champ des vitesses $\Vec{v}$.
% \item Si $\Vec{v}$ est porté par un vecteur de base $\Vec{u}$, les lignes de courant sont tangentes à $\Vec{u}$.
% \item Si $\Vec{v}$ a plusieurs composantes,
% \begin{liste} 
%   \item calculer et projeter l'équation $\Vec{v}\pv \Vec{\d l}=\Vec{0}$  ;
%   \item séparer les variables spatiales ($x,y$ en cartésiennes, $r,\theta$ en polaire, ...);
%   \item intégrer pour obtenir une équation de courbe.
% \end{liste}
% \item Tracer les lignes de courant en indiquant le sens de l'écoulement.
% \end{methode}
%
% \begin{liste}
%   \item Si $\Vec{v}(r,t)$ ne possède qu'une seule composante dans une base donnée, par exemple est porté par $\ex$ alors les lignes de courant sont données par l'équation \guill{coordonnée associée à la composante = constante}, soit ici $x=\cte$ : il s'agit d'une ligne droite horizontale.
%   \item Si $\Vec{v}(r,t)$ possède plusieurs composantes, les lignes de courant vérifient en coordonnées cartésiennes
%   \begin{equation}
%   \begin{pmatrix}
%   v_x \\ v_y \\ v_z
%   \end{pmatrix}
%   \pv
%   \begin{pmatrix}
%   \d x \\ \d y \\ \d z
%   \end{pmatrix}
%   = \begin{pmatrix}
%   v_y \d z - v_z \d y \\ v_z \d x - v_x \d z \\ v_x \d y - v_y \d x
%   \end{pmatrix}
%   =
%   \begin{pmatrix}
%   0\\0\\0
%   \end{pmatrix},
%   \end{equation}
%   et ont donc pour équation
%   \begin{equation}
%   \begin{split}
%   \dfrac{\mathrm{d}x}{v_{x}}=\dfrac{\mathrm{d}y}{v_{y}}=\dfrac{\mathrm{d}z}{v_{z}}.
%   \end{split}
%   \end{equation}
%   On peut trouver des formules équivalentes dans les autres systèmes de coordonnées.
% \end{liste}
%\begin{remarque}
%En coordonnées cylindriques et sphériques on obtient respectivement
%\begin{equation}
%\begin{split}
%\dfrac{\mathrm{d}r}{v_{r}}=\dfrac{r\,\mathrm{d}\theta}{v_{\theta}}=\dfrac{\mathrm{d}z}{v_{z}}\qquad\text{et}\qquad \dfrac{\mathrm{d}r}{v_{r}}=\dfrac{r\,\mathrm{d}\theta}{v_{\theta}}=\dfrac{r\sin\theta\,\mathrm{d}\varphi}{v_{\varphi}}.
%\end{split}
%\end{equation}
%\end{remarque}

\begin{fig}[flow]
  \centering
  \includegraphics[width=0.7\columnwidth]{flow.png}
  \capt{Visualisation des lignes de courant (de gauche à droite) dans un écoulement stationnaire autour d'une d'aile d'avion (au centre de la figure). Les lignes de courant sont obtenues en injectant du colorant en certains points de l'écoulement. Crédits : ONERA Photograph, Werlé 1974.}
\end{fig}

% \begin{exemple}{Détermination d'une ligne de courant}
%   On considère le champ des vitesses défini dans un repère cartésien par $\Vec{v}(x,y,t) = -kx \ex + ky \ey,$
%   avec $k$ une constante. Déterminer les lignes de courant.
% \end{exemple}
% \examplesol{
% La vitesse a deux composantes $v_x = -kx$ et $v_y = ky$.
% \begin{liste}
%   \item on calcule   $$
%     \begin{pmatrix}
%     v_x \\ v_y \\ 0
%     \end{pmatrix}
%     \pv
%     \begin{pmatrix}
%     \d x \\ \d y \\ \d z
%     \end{pmatrix}
%     = \begin{pmatrix}
%     v_y \d z \\  - v_x \d z \\ v_x \d y - v_y \d x
%     \end{pmatrix}
%     =
%     \begin{pmatrix}
%     0\\0\\0
%     \end{pmatrix},
%     $$
%     ce qui donne $\d z = 0$ soit $z=\cte$ (les lignes de courant sont planes) et $v_x \d y = v_y \d x$.

%     \item Séparation des variables, $ - k x \d y = k y \d x \Rar -\frac{\d x}{x} = \frac{\d y}{y}$.
%     \item Intégration  $-\ln(x) = \ln y + {\rm Cte} \Rightarrow y(x) = \frac{a}{x},$
%     avec $a$ une constante. Il s'agit d'une branche d'hyperbole.
% \end{liste}
% \begin{fig}[courant]
%   \centering
%   \includegraphics[scale=1]{ligne_courant.pdf}
%   % \capt{Exemple d'écoulement}
% \end{fig}
% % \end{enumerate}
% }

On peut visualiser ces lignes expérimentalement en suivant le déplacement de traceurs dans l'écoulement. Enfin, on appelle \imp{tube de courant} un ensemble de lignes de courant qui s'appuient sur un contour fermé. Un tube de courant est donc une canalisation, réelle ou fictive, dans l'écoulement.

\subsection{Notion de dérivée particulaire}

\parag{Dérivée particulaire}
En description eulérienne des fluides, les champs dépendent de l'espace et du temps. On souhaite suivre l'évolution dans le temps d'une propriété d'une particule fluide $\Pf$, par exemple, sa masse volumique $\rho_{\Pf}$. À l'instant $t$, la particule de fluide occupe un point $\rm M$ de vecteur position $\Vec{r}$ et possède une température $\rho_{\Pf}(t)=\rho(\Vec{r},t)$. À un instant voisin $t+\d t$, cette même particule occupe la position $\Vec{r} + \d \Vec{r}$ et possède une température $\rho_{\Pf}(t+\d t)=\rho(\Vec{r}+\Vec{\d r},t+\d t)$. Le déplacement de la particule est provoqué par l'écoulement : le déplacement élémentaire s'écrit donc $\Vec{\d r} = \Vec{v}(M,t) \d t$, où $\Vec{v}(M,t)$ est le champ des vitesses.
\begin{figure}[H]
  \centering
  \includegraphics{derivee_particulaire.pdf}
  \caption{Évolution dans le temps d'une particule de fluide $\Pf$ entre deux instants voisins $t$ et $t+\d t$. La particule de fluide est caractérisée par sa masse volumique $\rho_{\Pf}(t)$.}
  \label{fig:}
\end{figure} 

Calculons maintenant le taux de variation de masse volumique de la particule fluide :
\begin{align}
  \dd{\rho_{\Pf}}{t}(t) = 
\dfrac{\rho_{\Pf}(t+\d t)-\rho_{\Pf}(t)}{\d t} & = \dfrac{\rho(\Vec{r}+\Vec{\d r},t+\d t) -(\rho\Vec{r},t)}{\d t} \\ & = \underbrace{\dfrac{\rho(\Vec{r},\JRouge{t+\d t})-T(\Vec{r},\JRouge{t})}{\d t}}_{\JRouge{\fbox{1}}} + \underbrace{\dfrac{\rho(\JVert{\Vec{r}+\Vec{\d r}},t+\d t)-T(\JVert{\Vec{r}},t+\d t)}{\d t}}_{\JVert{\fbox{2}}}.
\label{eq:varia_T}
\end{align}
La variation de température de l'\Eq{varia_T} a donc deux contributions.
\begin{enumerate}
\item Le premier terme provient de l'évolution temporelle du champ de température en un point $\rm M$ donné, et tend vers $\dr{\rho}{t}(\Vec{r},t)$ lorsque $\d t$ tend vers $0$. On appelle dérivée \imp{locale} l'opérateur 
\eqbox[]{
\dr{}{t}.
}
\item Le second terme provient de l'évolution spatiale du champ de température en un instant $t$ donné. Effectuons un développement de Taylor lorsque $\d t$ tend vers $0$ en utilisant les coordonnées cartésiennes  :
\begin{align*}
\rho(\Vec{r}+\Vec{\d r},t+\d t) - \rho(\Vec{r},t+\d t) & =\rho(x +\d x, y + \d y,z + \d z, \d t,t+\d t) - \rho(x,y,z,t+\d t) \\
& = \d x \dr{\rho}{x}(\Vec{r},t+\d t) + \d y \dr{\rho}{y}(\Vec{r},t+\d t) + \d z \dr{\rho}{z}(\Vec{r},t+\d t) \\ & = \Vec{\d r} \cdot (\Grad{\rho})(\Vec{r},t+\d t) \\
& =  \Vec{v} \d t \cdot (\Grad{\rho})(\Vec{r},t+\d t).
\end{align*}
On a donc 
\begin{equation}
\dfrac{\rho(\Vec{r}+\Vec{\d r},t ) -\rho(\Vec{r},t)}{\d t} = (\Vec{v} \cdot \Grad{T})(\Vec{r},t).
\end{equation}
On appelle dérivée \imp{convective} l'opérateur 
\eqbox[]{
\Vec{v} \cdot \Grad{}.
}
\end{enumerate}

\begin{definition}[Dérivée particulaire]\label{def:D_part}
On appelle \imp{dérivée particulaire} d'un champ scalaire $f(M,t)$ le taux de variation de la grandeur $f$ en suivant la particule de fluide dans son déplacement, d'expression
\begin{equation*}
\Dt{f} = \dr{f}{t} + (\Vec{v} \cdot \grad){f}.
\end{equation*}
\end{definition}


\parag{Accélération particulaire} La dérivée particulaire s'applique aussi à un champ vectoriel, comme la vitesse $\Vec{v}$. De façon naturelle, la dérivée particulaire du champ de vitesse s'appelle l'accélération particulaire, notée $\Vec{a}(M,t)$, et représente l'accélération d'une particule de fluide en $M$ à l'instant~$t$,
% \begin{equation}
\eqbox{
  \Vec{a} = \dfrac{\mathrm{D}\Vec{v}}{\mathrm{D}t} =\dfrac{\partial \Vec{v}}{\partial t}+(\Vec{v}\cdot\Grad)\Vec{v}.
  }
% \end{equation}
\begin{rappel}{Opérateur de dérivée convective}
On rappelle que l'opérateur gradient correspond formellement au vecteur $\left(\dr{}{x},\dr{}{y},\dr{}{z}\right)$. Cet opérateur ne peut s'appliquer qu'à un champ scalaire ! En coordonnées cartésiennes, on a par exemple
\begin{equation}
\Grad{T} = \left(\dr{T}{x},\dr{T}{y},\dr{T}{z}\right).
\end{equation}
La dérivée convective peut s'appliquer à un champ scalaire ou vectoriel. En coordonnées cartésiennes, elle s'exprime comme
$
\Vec{v}\cdot\Grad{} =
v_{x}\dfrac{\partial}{\partial x} + v_{y}\dfrac{\partial}{\partial y} + v_{z}\dfrac{\partial}{\partial z}
$. Appliqué au champ de vitesse, cela donne
\begin{equation}
(\Vec{v}\cdot\Grad)\Vec{v}=
\left(\begin{array}{ccc}
v_{x}\dfrac{\partial \JRouge{v_{x}}}{\partial x}+v_{y}\dfrac{\partial \JRouge{v_{x}}}{\partial y}+v_{z}\dfrac{\partial \JRouge{v_{x}}}{\partial z} \\ [3mm]
v_{x}\dfrac{\partial \JRouge{v_{y}}}{\partial x}+v_{y}\dfrac{\partial \JRouge{v_{y}}}{\partial y}+v_{z}\dfrac{\partial \JRouge{v_{y}}}{\partial z} \\ [3mm]
v_{x}\dfrac{\partial \JRouge{v_{z}}}{\partial x}+v_{y}\dfrac{\partial \JRouge{v_{z}}}{\partial y}+v_{z}\dfrac{\partial \JRouge{v_{z}}}{\partial z} \\  [3mm]
\end{array}\right).
\end{equation}
Dans les autres systèmes de coordonnées, soit l'expression vous sera donné, soit il faudra utiliser une autre formule d'analyse vectorielle (voir fin du chapitre).
\end{rappel}


\section{Conservation de la masse}
\subsection{Débits volumique et massique}

Un écoulement est un phénomène de transport.  Cherchons à suivre le volume et la masse d’une particule de fluide emportée par l’écoulement.


\parag{Débit volumique} Soit une section élémentaire de normale $\dS = \d S \Vec{n}$. Pendant une durée élémentaire $\d t$, le fluide au voisinage du point $M$ se déplace d’une quantité $\Vec{\d r} = \Vec{v}(M,t)\,\d t$. Le volume de fluide qui traverse la surface élémentaire pendant $\d t$ est donc donné par (\Fig{Dv})
\begin{equation}
\dV = \d S \times \Vec{\d r}\cdot\Vec{n} = \Vec{v}(M,t)\,\d t\cdot\dS.
\end{equation}
%%%%%%%%%%%%%%%%%%%%%%%%%
\begin{figure}[H]
\centering
\includegraphics[scale=1]{Qm.pdf}
\caption{À partir d’une surface $S$, on peut définir un tube de courant élémentaire en suivant la trajectoire du contour pendant un temps élémentaire $\d t$.}
\label{fig:Dv}
\end{figure}
%%%%%%%%%%%%%
%\begin{remarque}
%La notation $\d^{2} V$ souligne que la quantité est infinitésimale du second ordre, car infinitésimale en espace (facteur $\d S$) \imp{et} en temps (facteur $\d t$). En quelque sorte, $\d^2 V=\d(\d V)$ : cette notation est ainsi à rapprocher de celle qui intervient au numérateur lorsqu'on note la dérivée seconde de la position (accélération) comme $\frac{\d^2 x}{\d t^2}$.
%\end{remarque}
On appelle débit volumique élémentaire $\d \DV$ le volume élémentaire de fluide qui traverse la surface élémentaire $\Vec{\d S}$ par unité de temps, autrement dit
\begin{equation}\label{eq:dDV}
\d\DV=\frac{\dV}{\d t}=\Vec{v}(M,t)\cdot \dS.
\end{equation}
On obtient ensuite le débit volumique total en intégrant l’\Eq{dDV} sur une surface macroscopique $S$.
\begin{definition}
On appelle \imp{débit volumique} le volume traversant une section $\mc{S}$ par unité de temps, d’expression
\begin{equation*}
\DV  = \dd{V}{t}= \iint_{\mc{S}}\Vec{v}\cdot \dS.
\end{equation*}
%Le débit volumique s’exprime en $\SI{}{m^2.s^{-1}}$.
\end{definition}
\noindent Le débit volumique s'identifie ainsi au flux du vecteur vitesse. Il s'exprime en $\SI{}{m^3.s^{-1}}$ et dépend en général du temps $t$ et de la section $S$.
%\begin{remarque}
%Si l'on introduit la quantité
%\begin{equation}
%\d V=\iint_S \d^{2} V =\iint_{S}\Vec{v}(M,t)\,\d t\cdot\Vec{\d S},
%\end{equation}
%alors on peut écrire que
%\begin{equation}
%D_V=\frac{\d V}{\d t}=\iint_{S}\Vec{v}(M,t)\cdot\Vec{\d S}.
%\end{equation}
%Cette expression est par exemple à rapprocher des expressions du flux dans les chapitres de diffusion de particules et de diffusion thermique.
%\end{remarque}


\parag{Débit massique} De même, la masse élémentaire traversant la surface élémentaire  $\dS$ pendant un temps $\d t$ s'écrit
\begin{equation}\label{eq:dDV}
\dm = \rho(M,t) \dV = \rho(M,t)\Vec{v}(M,t)\,\d t\cdot\dS,
\end{equation}
où $\rho$ désigne la masse volumique. Intégrons cette relation sur une section macroscopique $\mc{S}$.
\begin{definition}\label{def:Dm}
On appelle débit massique la masse traversant une section $\mc{S}$ par unité de temps, d’expression
\begin{equation*}
\Dm = \frac{\d m}{\d t}=\iint_{\mc{S}}\rho\Vec{v}\cdot \dS.
\end{equation*}
\end{definition}
\noindent Le débit massique s'exprime en $\SI{}{\kg.s^{-1}}$ et dépend en général du temps $t$. 
En général, la masse volumique $\rho(M,t)$ dépend du point $M$ : on ne peut pas la sortir de l'intégrale de la \Def{Dm}. En revanche, dans un fluide homogène, $\rho$ ne dépend pas de l'espace et donc $\Dm = \rho \DV$.

\begin{remarque}
  Le débit massique s'identifie au flux du vecteur \imp{densité de courant de masse}
  \begin{equation}
  \Vec{\jmath}_m=\rho\Vec{v}.
  \end{equation}
  Noter l’analogie avec la notion de flux thermique ou de flux de particules :
  \begin{equation}
    \Pth = \iint_S \Vec{\jth} \cdot \dS \quad \sim \quad \Phi_N = \iint_S \Vec{\jN} \cdot \dS \quad \sim \quad   \DV = \iint_S \Vec{v} \cdot \dS \quad \sim \quad \Dm = \iint_S \Vec{\jmath}_m \cdot \dS.
  \end{equation}
\end{remarque}



\subsection{Équation de conservation de la masse}

La masse est une grandeur conservative (en l'absence de réaction nucléaires) : la masse d'une particule de fluide emportée par l'écoulement est constante. De façon équivalente, la variation de la masse contenue dans un volume fixe ne peut provenir que du transport de matière par l'écoulement à travers le volume. %Ces échanges sont quantifiés par le débit massique.
Traduisons cette loi de conservation sur un exemple uni-dimensionnel. La méthode est très générale et doit être maîtrisée.


\begin{methode}[Établir l'équation de conservation de la masse]
  \item Définir le système comme une portion élémentaire du milieu :
  \begin{liste}
    \item tronçon d'épaisseur $\d x$ et de section $S$ en cartésienne ;
    \item coquille cylindrique de rayon $r$, d'épaisseur $\d r$ et de hauteur $H$ en cylindrique ;
    \item coquille sphérique de rayon $r$, d'épaisseur $\d r$ en sphérique.
  \end{liste}
  \item Exprimer la variation de masse $\d m$ pendant $\d t$ en fonction de $\dr{\rho}{t}$.
  \item Exprimer cette variation comme la masse échangée $\d m = \delta m_{e} - \delta m_s$, puis en fonction de $\rho \Vec{v}$.
  \item Écrire l’égalité entre ces deux expressions et simplifier.
\end{methode}


\parag{Équation locale} Appliquons la méthode ci-dessus.

\begin{enumerate}
  \item Considérons un tube de courant élémentaire dirigé selon $x$, de section $S$ et situé entre les abscisses $x$ et $x + \d x$. L'écoulement est localement dirigé suivant $x$. Pour simplifier, supposons que la vitesse ne dépend que de l'abscisse $x$ et du temps $t$, de sorte que $\Vec{v}=v_{x}(x,t)\ex$. Réalisons un bilan de masse sur ce tube de courant, de volume $\d V=S\d x$.
  \item  Pendant un temps $\d t$, la variation de masse à l'intérieur du système est égale à la différence de ce qui est présent à l'instant final $t+\d t$ moins ce qui était présent à l'instant initial $t$, d'où
  \begin{equation}
  \begin{split}
  \d m&=m(x,t+\d t)-m(x,t)=\rho(x,t+\d t)S\d x-\rho(x,t)S\d x=\frac{\partial \rho}{\partial t} \d t S\d x.
  \end{split}
  \end{equation}
  \item 
  Cette variation de masse provient de ce qui est entré dans le tube de courant à l'abscisse $x$, moins ce qui est sorti à l'abscisse $x+\d x$ pendant ce même temps $\d t$, c'est-à-dire
  \begin{equation}
  \begin{split}
  \d m=\delta m_{\rm entrée}-\delta m_{\rm sortie}=\rho(x,t)v_{x}(x,t)S\d t - \rho(x,t)v_{x}(x+\d x,t)S\d t=-\frac{\partial \rho v_x}{\partial x}\d x S\d t.
  \end{split}
  \end{equation}
  \item 
  En combinant les deux expressions, on obtient
  \begin{equation}\label{eq:conti_1d}
  \dfrac{\partial \rho}{\partial t}+\frac{\partial \rho v_x}{\partial x}=0.
  \end{equation}
  
\end{enumerate}

%\begin{rappel}{Opérateur divergence}
%En généralisant l'\Eq{conti_1d} à trois dimensions, on obtiendrait
%\begin{equation}
%\dfrac{\partial \rho}{\partial t}+\frac{\partial j_x}{\partial x} + \frac{\partial j_y}{\partial y} + \frac{\partial j_z}{\partial z} =0.
%\end{equation}
%On introduit l'opérateur divergence, noté $\Dive{}$, tel que pour tout vecteur $\Vec{u} = (u_x,u_y,u_z)$,
%\begin{equation}
%\Dive{\Vec{u}} = \dr{u_{\JRouge{x}}}{\JRouge{x}} + \dr{u_{\JRouge{y}}}{\JRouge{y}} + \dr{u_{\JRouge{z}}}{\JRouge{z}}.
%\end{equation}
%L'opérateur divergence prend un vecteur et renvoie un scalaire.
%\end{rappel}
On obtient l'équation de conservation de la masse, aussi appelée équation de continuité. À trois dimensions, on remplace l'opérateur $\dr{}{x}$ par l'opérateur $\dive{}$.
 % après remplacement de $\Vec{j} = \rho \Vec{v}$.
\begin{theorem}[Équation de conservation de la masse]\label{form:conserv_m}
\begin{equation*}
\dfrac{\partial \rho}{\partial t}+\Dive{(\rho \Vec{v})}=0.
\end{equation*}
\end{theorem}
\noindent Il s'agit d'une équation locale, puisqu'elle relie la vitesse $\Vec{v}(M,t)$ et la masse volumique $\rho(M,t)$ en un point $M$ et un instant $t$ donnés.

%\begin{rappel}{Démonstration}
%Voici une autre démonstration possible, utilisant l'analyse vectorielle.
%\begin{equation}
%\frac{\partial m}{\partial t} =\frac{\partial}{\partial t}\iiint_{V}\rho\,\mathrm{d}V\\
%= D_{m_{\rm{e}}} - D_{m_{\rm{s}}}=-\int_{S}\Vec{j}\cdot \Vec{\mathrm{d}S}.
%\end{equation}
%Or d'après le théorème de Green-Ostrogradski,
%\begin{equation}
%\int_{S}\Vec{j}\cdot \Vec{\mathrm{d}S}=\iiint_{V}(\Dive{\Vec{j}})\,\mathrm{d}V.
%\end{equation}
%Ainsi le bilan de masse se simplifie en
%\begin{equation}
%\iiint_{V}\left(\dfrac{\partial \rho}{\partial t}+\Dive{\Vec{j}}\right)\mathrm{d}V=0
%\end{equation}
%Cette relation étant vraie quel que soit $\d V$, l'intégrande est nulle, d'où
%\begin{equation}
%\dfrac{\partial \rho}{\partial t}+\Dive{\Vec{j}}=0.
%\end{equation}
%\end{rappel}


\begin{remarque}
  \udl{Autre écriture}.
% \parag{Interprétation physique}
% \label{subsec:alternative_cons_masse}
En utilisant l'identité vectorielle $\Dive{(\rho\Vec{v})}=\rho\,\Dive{\Vec{v}}+\Vec{v}\cdot\Grad{\rho}$, on peut récrire l'équation de conservation de la masse sous la forme
\begin{align}\label{eq:alternative_conserv_masse}
0  = \dfrac{\partial \rho}{\partial t}+\Dive{(\rho \Vec{v})}  = \JRouge{\frac{\partial \rho}{\partial t}+\Vec{v}\cdot\Grad{\rho}}+\rho\,\Dive{\Vec{v}} = \JRouge{\frac{\mathrm{D}\rho}{\mathrm{D} t}}+\rho\,\Dive{\Vec{v}},
\end{align}
en reconnaissant l'expression de la dérivée particulaire de $\rho$.
% Cette formulation nous permet d'interpréter la divergence du champ de vitesse.
\end{remarque}

\begin{exemple}{Interprétation de la divergence du champ des vitesses}
  Nous représentons ci-dessous l'allure du champ des vitesses en fonction du signe de $\dive{\Vec{v}}$, comme indiqué dans le  \refFormu{B}.
  \begin{fig}[label]
    \centering
    % \subfigure[]{
    \def\WID{0.2\columnwidth}
\includegraphics[height=\WID]{crop_divergence1.pdf}
% }
% \subfigure[]{
\includegraphics[height=\WID]{crop_divergence2.pdf}
% }
% \subfigure[]{
\includegraphics[height=\WID]{crop_divergence3.pdf}
% }
% \subfigure[]{
\includegraphics[height=\WID]{crop_divergence4.pdf}
% }
    %\capt{}
  \end{fig}
  \begin{quest}
    \item	À partir de l'équation de conservation de la masse, déterminer le signe de $\dive{\Vec{v}}$ en fonction de la variation temporelle de la masse volumique $\rho_\Pf(t)$ d'une particule de fluide.
    \item En déduire la variation du volume de la particule de fluide dans chacune des situations. Représenter graphiquement l'évolution de la particule de fluide, et comparer aux schémas ci-dessus.
  \end{quest}
\end{exemple}
\examplesol{
  De façon générale, le signe de $\dive \Vec{v}$ donne une information sur l'évolution du volume des particules de fluide. En effet, d'après  l'\Eq{alternative_conserv_masse},
	\begin{liste}
	  \item Si $\dive \Vec{v} < 0$, $\Dt{\rho} > 0$ donc $\rho$ augmente et le volume diminue : la particule de fluide se contracte ;
	  \item Si $\dive \Vec{v} = 0$,  $\Dt{\rho} = 0$ donc $\rho$ est constant : la particule de fluide est de volume constant ;
	  \item Si $\dive \Vec{v} > 0$, $\Dt{\rho} < 0$ donc $\rho$ diminue et le volume augmente : la particule de fluide se dilate.
	\end{liste}
}

% En effet, $\Dive{\Vec{v}}=-\frac{1}{\rho}\frac{\mathrm{D}\rho}{\mathrm{D}t}$. Or le produit $\rho V=m$ représente la masse constante d'une particule de fluide. Par passage au logarithme, $\ln(m)=\ln(\rho)+\ln(V)$, et donc en dérivant,
% $
% \Dt{\rho}+\Dt{V}=0.
% $
% Autrement dit,
% \begin{equation}
% \label{eq:interp_div_variation_volume}
% \Dive{\Vec{v}}=\frac{1}{V}\frac{\mathrm{D}V}{\mathrm{D}t}.
% \end{equation}
% La divergence du champ de vitesse indique la variation relative de volume d'une particule de fluide par unité de temps. Dans un écoulement divergent, une particule de fluide se déforme au cours de son mouvement. Lorsque  $\Dive{\Vec{v}}>0$, la particule de fluide s'agrandit ; lorsque  $\Dive{\Vec{v}}<0$, elle se rétrécit.
%Même si le vecteur vitesse de l'écoulement dépend du référentiel, la divergence du vecteur vitesse est indépendante du référentiel. C'est ce que montre la démonstration précédente, puisque $\Dive{\Vec{v}}$ s'identifie à la variation relative de volume de la particule de fluide\footnote{Pour une étude détaillée, on se reportera au Complément~\ref{comp:déformation_carré}.}, qui elle ne dépend pas du référentiel.
% Cette propriété indique par ailleurs que le caractère divergent d'un écoulement ne dépend donc pas du référentiel.
% Pour quelques animations d'écoulements classiques, voir par exemple \href{https://femto-physique.fr/mecanique_des_fluides/cinematique.php}{\JRouge{ici}}.

\section{Écoulements particuliers}
\subsection{L'écoulement stationnaire}
%« Stationnaire » signifie, indépendants du temps. Dans le cadre de la mécanique des fluides, un écoulement est dit stata

\begin{definition}\label{def:stat}
Un écoulement est \imp{stationnaire}, ou permanent, lorsque tous les champs eulériens : vitesse $\Vec{v}(M,t)$, masse volumique~$\rho(M,t)$,  etc., sont indépendants du temps $t$. Formellement, on peut écrire
$$\frac{\partial }{\partial t} = 0.$$
\end{definition}
%Tous les champs eulériens ($P,v,\rho ...$) sont indépendants du temps,.
 %
%Autrement dit, toutes les particules de fluide qui passent en un point donné de l'écoulement ont \textit{à tout instant} le même vecteur vitesse. Uniquement dans ce cas,
% Dans ce cas uniquement, les lignes de courant, que l'on peut observer sur une carte de champ, s'identifient aux trajectoires des particules de fluide.

\begin{remarque}
Attention, le caractère stationnaire d'un écoulement dépend du référentiel choisi ! L'écoulement d'une rivière est stationnaire pour une barque emportée par l'eau, auquel cas la rivière apparaît même statique, mais pas pour un pêcheur situé sur la rive.
\end{remarque}

\parag{Conservation du débit volumique}
Dans un écoulement stationnaire, on a en particulier $\dr{\rho}{t} = 0$. L'équation de conservation de la masse \Formu{conserv_m} se simplifie en
\begin{equation}
\dfrac{\partial \rho}{\partial t}+\Dive{(\rho \Vec{v})}=\Dive{(\rho \Vec{v})}=0.
\end{equation}

\begin{rappel}{Théorème de Green-Ostrogradski}
Ce théorème d'analyse vectorielle affirme que pour tout vecteur $\Vec{j}$,
%permet de relier l'intégrale de la divergence d'un vecteur $\Vec{j}$ sur un volume $V$ au flux de $\Vec{j}$ sur la surface fermée $S$ qui délimite le volume $V$,
%\begin{equation}
%\frac{\partial m}{\partial t} =\frac{\partial}{\partial t}\iiint_{V}\rho\,\mathrm{d}V\\
%= D_{m_{\rm{e}}} - D_{m_{\rm{s}}}=-\int_{S}\Vec{j}\cdot \Vec{\mathrm{d}S}.
%\end{equation}


\begin{equation}
\iiint_{V}(\Dive{\Vec{j}})\,\mathrm{d}V = \oiint_{S}\Vec{j}\cdot \Vec{\mathrm{d}S}.
\end{equation}
où $S$ est la surface fermée qui délimite le volume $V$. Le rond qui entoure la double intégrale est là pour préciser que l'ensemble d'intégration (ici la surface $S$) est fermé. Le vecteur $\dS$ est orientée vers l'extérieur du volume entouré.

\end{rappel}

Appliquons le théorème de Green-Ostrogradski sur un tube de courant délimité par les surfaces entrante $S_{1}$, sortante $S_{2}$ et latérale $S_{\rm lat}$, \Fig{conduite_incompr}. La surface $S = S_{1} \cup S_{\rm lat} \cup S_{2}$ est fermée, et donc
\begin{equation}
  \iiint_{V}\Dive{(\rho\Vec{v})} \,\d V = \oiint_{S}\rho \Vec{v}\cdot \dS = \iint_{S_1}\rho \Vec{v}\cdot \dS + \iint_{S_{\rm lat}}\rho \Vec{v}\cdot \dS + \iint_{S_2}\rho \Vec{v}\cdot \dS = \Dm^{(\rm 2)} - \Dm^{(\rm 1)},
\end{equation}
puisque l'intégrale sur la surface latérale est nulle, et où $\Dm$ est défini positif dans le sens de l'écoulement par convention. Pour un écoulement stationnaire, on a donc $\Dm^{(\rm 2)} = \Dm^{(\rm 1)}$.
 % Le débit massique sortant $\Dm^{(\rm s)}$ égale le débit massique entrant $\Dm^{(\rm e)}$.

\begin{prop}\label{prop:Dm}
Dans un écoulement stationnaire, le \imp{débit massique} se conserve dans un tube de courant.
\end{prop}
% \noindent Cette relation traduit en fait quelque chose de bien connu : la loi des nœuds (\Fig{conserv_débit}) !

\begin{exemple}{Conservation du débit massique}
  Appliquer la conservation du débit massique aux situations suivantes.
  \setcounter{quest}{0}
\begin{quest}
  \item Vitesse uniforme sur toute section $S$ d'un tube de courant.
  \item Conduite de section $S_0$ qui se scinde en deux petites conduites de sections $S_1$ et $S_2$.
\end{quest}
\end{exemple}
\examplesol{
\setcounter{quest}{0}
\begin{quest}
  \item
  Le débit massique vaut $\Dm = \rho S v$ et est constant le long d'un tube de courant. On a donc $\rho_1 S_1 v_1=\rho_2S_2v_2$.
  \item
  On a $\Dm^{(0)} =  \Dm^{(1)}  + \Dm^{(2)}$, il s'agit de  la loi des nœuds hydraulique.
\end{quest}
}

\parag{Stationnarité approchée} La plupart des écoulements ne sont pas strictement stationnaires, mais peuvent l'être de façon approchée. À quelle condition ? Il suffit pour cela que les dérivées partielles par rapport au temps soient « petites ». En particulier, l'accélération locale doit être négligeable,
$$
\norm{\dr{\Vec{v}}{t}} \ll \norm{(\Vec{v}\ps\grad) \Vec{v}}.
$$
Si $U$, $T$ et $L$ sont les vitesse, temps et taille caractéristiques de l'écoulement, alors on doit avoir
$$
\norm{\dr{\Vec{v}}{t}} \sim \frac{U}{T} \ll \norm{(\Vec{v}\ps\grad) \Vec{v}} \sim \frac{U^2}{L} \Rar \boxed{U \gg \frac{L}{T}}.
$$

%Cette relation étant vraie quel que soit $\d V$, l'intégrande est nulle, d'où
%\begin{equation}
%\dfrac{\partial \rho}{\partial t}+\Dive{\Vec{j}}=0.
%\end{equation}

%D'après le théorème de Green-Ostrogradski, cette relation implique que
%\begin{equation}
%\begin{split}
%\oiint_{\mathcal{S}} \Vec{j}\cdot \Vec{\mathrm{d}S}=0\qquad\Leftrightarrow\quad \iint_{\mathcal{S}} \Vec{j}\cdot \Vec{\mathrm{d}S}=\cte~~(\text{indép. de }\mathcal{S}).
%\end{split}
%\end{equation}
%Les deux formulations sont équivalentes, dans la première on considère une surface fermée (d'où le rond qui entoure la double intégrale), alors que dans la deuxième la surface est quelconque et on dit simplement que le flux à travers cette surface ne dépend pas du choix de la surface.
%En conclusion, le flux de $\Vec{j}$ à travers une surface fermée (un tube de courant) est conservé: on dit que $\Vec{j}$ est à flux conservatif. Ainsi le long d'un tube de courant, le débit massique entrant $D_{m_{\rm{e}}}$ est égal au débit massique sortant $D_{m_{\rm{s}}}$. Donc le débit massique se conserve.



\subsection{L'écoulement incompressible}

Un écoulement est \imp{incompressible} si le volume d'une particule de fluide est constant. Cela signifie que la masse volumique $\rho(M,t)$ d'une particule de fluide est constante, soit d'après l'\Eq{alternative_conserv_masse},
$$
\Dt{\rho} = 0 \Rar 0 = \frac{\mathrm{D}\rho}{\mathrm{D} t} + \rho\,\Dive{\Vec{v}} = \rho \,\dive \Vec{v}.
$$


% \begin{align}\label{eq:alternative_conserv_masse}
% 0  = \dfrac{\partial \rho}{\partial t}+\Dive{(\rho \Vec{v})}  = \frac{\partial \rho}{\partial t}+\Vec{v}\cdot\Grad{\rho}+\rho\,\Dive{\Vec{v}} = \frac{\mathrm{D}\rho}{\mathrm{D} t}+\rho\,\Dive{\Vec{v}},
% \end{align}

\begin{definition}\label{def:icom}
Un écoulement est \imp{incompressible} si le volume d'une particule de fluide est constant, soit
\begin{equation*}
\Dive{\Vec{v}}=0.
\end{equation*}
\end{definition}

% Les remarques suivantes sont reléguées en TD
% De façon générale, le signe de $\dive \Vec{v}$ donne une information sur l'évolution du volume des particules de fluide. En effet, d'après  l'\Eq{alternative_conserv_masse},
% \begin{liste}
%   \item Si $\dive \Vec{v} < 0$, $\Dt{\rho} > 0$ donc $\rho$ augmente et le volume diminue : la particule de fluide se contracte ;
%   \item Si $\dive \Vec{v} = 0$,  $\Dt{\rho} = 0$ donc $\rho$ est constant : la particule de fluide est de volume constant ;
%   \item Si $\dive \Vec{v} > 0$, $\Dt{\rho} < 0$ donc $\rho$ diminue et le volume augmente : la particule de fluide se dilate.
% \end{liste}
%
% \begin{remarque}
%   Un \emph{fluide} est incompressible lorsque $\rho$ est une constante dans tout le fluide, alors qu'un \emph{écoulement} est incompressible lorsque $\rho$ est une constante en suivant une particule de fluide. Donc
%   %%\begin{danger}
%   %Attention aux interprétations avec les dérivées particulaires, car
%   \begin{equation}
%     \begin{cases}
%       \rho=\cte \Rightarrow \frac{\mathrm{D}\rho}{\mathrm{D}t}=0 &\quad\text{mais}\quad  \frac{\mathrm{D}\rho}{\mathrm{D}t}=0  \nRightarrow \rho=\cte. \\
%       \text{fluide incompressible} \Rightarrow \text{écoulement incompressible} &\quad\text{mais}\quad  \text{écoulement incompressible} \nRightarrow \text{fluide incompressible}.
%      \end{cases}
%   \end{equation}
% \end{remarque}
% \noindent La propriété ci-dessus découle de l'\Eq{interp_div_variation_volume}. En utilisant l'expression \eqref{eq:alternative_conserv_masse} de la conservation de la masse, un écoulement incompressible vérifie $\frac{\mathrm{D}\rho}{\mathrm{D}t}=0$ : les particules de fluide conservent leur volume, et donc leur masse volumique, au cours de l'écoulement. En conséquence, le caractère incompressible d'un écoulement ne dépend pas du référentiel.

% Ceci veut dire que si le fluide est incompressible alors l'écoulement est incompressible, mais que l'écoulement peut être incompressible sans que le fluide le soit ! La condition $\frac{\D\rho}{\D t}=0$ implique seulement que les particules de fluide conservent leur volume (donc leur masse volumique) au cours de l'écoulement, mais ne signifie pas que $\rho=\cte$ dans tout le fluide. Mathématiquement, $\rho=\cte\Rightarrow\frac{\D\rho}{\D t}=0$, mais la réciproque est fausse en général.
%%\end{danger}
\parag{Conservation du débit volumique}
Dans un écoulement incompressible, le champ de vitesse
$\Vec{v}$ est à flux conservatif. Or le flux de $\Vec{v}$ est le débit volumique $\DV$.
% D'après le théorème de Green-Ostrogradski, le débit volumique entrant dans un tube de courant  $\DV^{(\rm{e})}$  égale le débit volumique sortant $\DV^{(\rm{s})}$.
\begin{prop}\label{prop:Dm}
Dans un écoulement incompressible, le \imp{débit volumique} se conserve dans un tube de courant.
\end{prop}
\noindent La conservation du débit volumique conduit à une relation du type $S_1v_1=S_2v_2$ entre deux sections $S_1$ et $S_2$ d'un tube de courant, ce qui apporte une information intéressante pour résoudre un problème. En particulier, si $S_1>S_2$, alors $v_1<v_2$ : la vitesse augmente lorsque la section diminue (Fig.~\ref{fig:conduite_incompr}).
%\begin{figure}[h!]
%\centering
%\includegraphics[width=0.35\textwidth]{conduite.pdf}
%\caption{Illustration de la conservation du débit volumique sur un tube de courant pour un écoulement incompressible ($S_1>S_2$ donc $v_1<v_2$).}
%\label{fig:conduite_incompr}
%\end{figure}

\begin{figure}[t!]
\centering
\subfigure[]{\label{fig:conduite_incompr}
\includegraphics[scale=1.1]{conserva_QV.pdf}
}
% \hspace{2cm}
\subfigure[]{\label{fig:conserv_débit}
\includegraphics[scale=1]{noeuds_QV.pdf}
}
\caption{(a) Dans un écoulement stationnaire, le débit massique se conserve dans un tube de courant. Si les champs eulériens sont uniformes sur la section, alors $\rho_1 v_1  S_1 = \rho_2 v_2 S_2$. Dans un écoulement incompressible, le débit volumique se conserve dans un tube de courant. Si les champs eulériens sont uniformes sur la section, alors $ v_1 S_1 = v_2 S_2 $. En particulier, si $S_2 > S_1$, alors $v_2<v_1$ : le fluide ralenti dans une conduite qui s'élargit. (b) Si une conduite se divise en deux, on a $\Dm^{(0)} =  \Dm^{(1)}  + \Dm^{(2)}$. Cette propriété équivaut à la loi des nœuds en électrocinétique. }
\end{figure}





\parag{Incompressibilité approchée}

% Il se peut que dans un écoulement incompressible, la dérivée locale $\frac{\partial \rho}{\partial t}$ compense exactement la dérivée convective $\Vec{v}\cdot\Grad{\rho}$, sans qu'aucun de ces deux termes ne soit nul.
Dans la plupart des situations (sauf pour les écoulements très visqueux), on pourra considérer l'écoulement comme incompressible si
% \begin{equation}
\eqbox[]{
  \mathrm{Ma}^2\ll 1,
  }
% \end{equation}
où le nombre sans dimension ${\rm Ma} = \frac{v}{c_s}$ est appelé \imp{nombre de Mach}, et où $v$ désigne la vitesse typique de l'écoulement de fluide et $c_{\rm{s}}$ la vitesse du son dans le milieu. Dans l'air, $c_{\rm{s}}=\SI{340}{m.s^{-1}}$, et même pour des vents forts soufflant à $\SI{100}{km.h^{-1}} \simeq \SI{30}{m.s^{-1}}$, on trouve ${\rm Ma} = \SI{8e-3}{} \ll 1$.
Bien que l'air soit un fluide compressible, l'écoulement de l'air est incompressible en bonne approximation dans ces conditions.

%\subsection{L'écoulement tourbillonnaire}
%\label{subsec:ecoulement_tourbillonnaire}
%%Un exemple typique est l'écoulement engendré par une tornade, mais pas seulement! En principe, ce type d'écoulement peut apparaître dès que la variation du champ de vitesse a lieu dans une direction orthogonale à sa direction de propagation. Par exemple, dans le cas le plus simple, $\Vec{v}(M,t)=v_{x}(y)\ex$. On parle dans ce cas d'un écoulement de cisaillement, dont la carte de champ est représentée en Fig.~\ref{fig:cisaillement}.
%%
%%\begin{figure}[h!]
%%\centering
%%\sousfigure{cisaillement.pdf}{}{0.48}{cisaillement_2.pdf}{}{0.48}
%%\caption{Carte de champ d'un écoulement stationnaire de cisaillement, dont la vitesse est de la forme $\vec{v}=v_x(y)\ex=y\ex$. Une particule de fluide est représentée à l'instant $t$ sur la figure de gauche, et elle se déforme au cours du temps, comme on peut le voir sur la figure de droite. Ici, la particule de fluide décrit un mouvement de rotation autour de $\vec{e_z}$ à la vitesse angulaire donnée par la norme du vecteur tourbillon $\vec{\Omega}$.}
%%\label{fig:cisaillement}
%%\end{figure}
%
%Quelque soit l'écoulement, on définit le vecteur tourbillon, noté $\Vec{\Omega}$, tel que
%\begin{equation}
%\Vec{\Omega}=\frac{1}{2}\,\Rot{\Vec{v}}.
%\end{equation}
%Il représente le vecteur rotation instantané de la particule de fluide, autour de son centre. Pour comprendre d'où provient cette relation, on consultera le Complément~\ref{comp:déformation_carré}.
%%Dans l'exemple précédent, $\Rot{\Vec{v}}=\Vec{\nabla}\pv\Vec{v}=-\frac{\partial v_{x}(y)}{\partial y}\Vec{e_z}=-\Vec{e_z}$ donc la rotation est effectivement selon $-\Vec{e_z}$, c'est-à-dire dirigée vers le fond.
%
%\paragraph{Un exemple classique: le tube de vorticité de Rankine (tornade).}
%On considère un écoulement incompressible ($\Dive{\Vec{v}}=0$) et tourbillonnaire ($\Rot{\Vec{v}}=2\Vec{\Omega}\neq 0$), comme par exemple celui d'une tornade. On cherche à déterminer le champ des vitesses en tout point autour du cœur de la tornade, en supposant que
%\begin{equation}
%\Vec{\Omega} =
%\begin{cases}
%\Omega\,\Vec{e_z} & \mbox{si } r\leq R, \\
%0 & \mbox{si } r>R.
%\end{cases}
%\end{equation}
%L'invariance par rotation autour de $(Oz)$ et celle par translation selon $(Oz)$ conduit à supposer que $v(M)=v(r)$. De plus, par symétrie, $\Vec{v}$ est un vecteur orthoradial, dirigé par $\et$. Autrement dit, $\Vec{v}(M)=v(r)\et$. Pour calculer $v(r)$ sachant qu'on connaît $\Vec{\Omega}$, donc le rotationnel de $\Vec{v}$, on utilise le théorème de Stokes
%\begin{equation}
%  \oint_{\Gamma}\Vec{v}\cdot\Vec{\d l}=\iint_{S(\Gamma)}\Rot{\Vec{v}}\cdot\Vec{\d S}=\iint_{S(\Gamma)}2\Vec{\Omega}\cdot\Vec{\d S}.
%\end{equation}
%Puisqu'on dispose de deux expressions du vecteur tourbillon dans deux régimes distincts, on distingue deux cas.
%\begin{enumerate}
%\item Pour $r\leq R$, au vu de la géométrie du problème, on choisit un contour $\Gamma$ qui est un cercle de rayon $r\leq R$. La surface $S(\Gamma)$ qui s'appuie sur le contour circulaire est donc un cylindre. Comme $\Vec{\Omega}$ est dirigé selon $\Vec{e_z}$, seule la composante de la surface élémentaire en cylindrique dont la normale est selon $\Vec{e_z}$ nous intéresse: elle s'exprime comme $\Vec{\d S}=r\d r\d\theta \Vec{e_z}$. Comme la vitesse est selon $\et$, selon la composante orthoradiale du déplacement élémentaire en cylindrique nous intéresse: il s'agit de $\Vec{\d l}=r\d\theta\et$. Le théorème de Stokes s'écrit alors
%\begin{equation}
%\begin{split}
%  &\int_{0}^{2\pi}v(r)\,r\d\theta=\int_{0}^{2\pi}\int_{0}^{r}2\Omega\,r\d r\,\d\theta\\
%&\Leftrightarrow 2\pi r v(r)=2\pi\,2\Omega\left(\tfrac{r^2}{2}\right).
%  \end{split}
%\end{equation}
%et donc finalement
%\begin{equation}
%\Vec{v}(r\leq R)=r\Omega\et.
%\end{equation}
%\item Pour $r>R$, l'idée est la même avec un contour de rayon $r>R$. Cette fois le théorème de Stokes donne (bien noter qu'on coupe l'intégrale surfacique au rayon $r=R$ car au delà le vecteur tourbillon est nul!)
%\begin{equation}
%\begin{split}
%  &\int_{0}^{2\pi}v(r)\,r\d\theta=\int_{0}^{2\pi}\int_{0}^{R}2\Omega\,r\d r\,\d\theta\\
%&\Leftrightarrow 2\pi r v(r)=2\pi\,2\Omega\left(\tfrac{R^2}{2}\right).
%  \end{split}
%\end{equation}
%et donc finalement
%\begin{equation}
%\Vec{v}(r>R)=\frac{\Omega R^2}{r}\et.
%\end{equation}
%\end{enumerate}
%\begin{figure}[h!]
%\centering
%\includegraphics[width=0.3\textwidth]{tornade.pdf}
%\caption{Champ des vitesses au sein d'une tornade (modèle simple).}
%\end{figure}

\subsection{L'écoulement irrotationnel}

\parag{Opérateur rotationnel} Le rotationnel agit sur un champ de vecteurs $\Vec{v}(M,t)$ et renvoie un champ de vecteurs $\Rot{\Vec{v}}(M,t)$ tel que
% \begin{equation}
\eqbox{
  \Rot{\Vec{v}} = \left( \dr{v_z}{y} - \dr{v_y}{z} \right) \ex + \left( \dr{v_x}{z} - \dr{v_z}{x} \right)\ey  + \left( \dr{v_x}{y} - \dr{v_y}{x} \right)\ez.
  }
% \end{equation}
Dans les autres systèmes de coordonnées, il faut utiliser un formulaire pour exprimer $\rot \Vec{v}$.

\begin{remarque}
  Formellement, le rotationnel correspond au produit vectoriel entre l'opérateur \guill{nabla} $\Vec{\nabla} = (\dr{}{x},\dr{}{y},\dr{}{z})$ et le champ vectoriel~$\Vec{v}$. On peut donc écrire
\begin{equation}
\Rot{\Vec{v}} = \Vec{\nabla} \pv \Vec{v} = \begin{pmatrix}\dr{}{x} \\[2.5mm] \dr{}{y}\\[2.5mm] \dr{}{z}\end{pmatrix} \pv \begin{pmatrix}\displaystyle v_x \\[2.5mm] v_y\\[2.5mm]v_z\end{pmatrix} = \begin{pmatrix}\dr{v_z}{y} - \dr{v_y}{z} \\[2.5mm] \dr{v_x}{z} - \dr{v_z}{x}\\[2.5mm] \dr{v_x}{y} - \dr{v_y}{x}\end{pmatrix}.
\end{equation}
\end{remarque}
\begin{exemple}{Interprétation de l'opérateur rotationnel}
  Étudions l'évolution d'une particule de fluide en rotation à la vitesse angulaire $\Omega$ autour de l'axe $\ez$. L'écoulement bidimensionnel associé est de champ des vitesses (en coordonnées cartésiennes)
  \begin{equation}
    \Vec{v}(x,y) = -\Omega y \,\ex + \Omega x\, \ey.
  \end{equation}
  \begin{fig}
    \centering
    \includegraphics[scale=1]{deformations_carre_2.pdf}
    \capt{Rotation d'une particule de fluide durant $\d t$. La particule, modélisée par un carré, tourne d'un angle $\d \alpha = \Omega \d t$, où $\Omega$ désigne la vitesse angulaire de rotation.}
  \end{fig}
  \begin{quest}
    \item	Justifier que ce champ des vitesses décrit bien la rotation en bloc d'une particule de fluide, en comparant la position du segment $\Vec{AB}$ aux instants $t$ et $t+\d t$.
    \item Montrer que $\rot \Vec{v} = 2\Omega \ez$. En déduire une interprétation du rotationnel.
  \end{quest}


    % Nous allons étudier le rôle de chaque terme entre crochets séparément.

\end{exemple}
\examplesol{
  On modélise une particule de fluide par un carré infinitésimal de côtés $\d x$ et $\d y$, et de sommets A, B, C et D, où A=O est l'origine du repère. Soit $\Vec{AB}=\d x \ex $ un segment à l'instant $t$ et  $\Vec{A'B'}$  ce même segment à l'instant $t+\d t$. On a
  \begin{equation}
    \begin{split}
    \label{eq:comp_deformation_generale}
    \d \Vec{AB} = \Vec{A'B'}-\Vec{AB}&= \Vec{BB'} - \Vec{AA'} =\left[\Vec{v}(B)-\Vec{v}(A)\right]\d t = \Vec{v}(B) \d t =  (\Omega \d x   \, \ey) \d t =  (\Omega \ez \pv \Vec{AB}) \d t,
    \end{split}
    \end{equation}
et donc
\begin{equation}
 \dd{\Vec{AB}}{t} = \Omega \ez \pv \Vec{AB}.
\end{equation}
On pourrait montrer que les autres segments du carré obéissent à la même équation. Cela montre que la particule de fluide tourne en bloc à la vitesse angulaire $\Omega$ autour de l'axe $(Oz)$ (dans le plan). On voit qu'on peut obtenir le vecteur rotation $\Omega \ez$ avec la relation
\begin{equation}
\Omega \ez = \frac{1}{2}\left( \dr{v_y}{x} - \dr{v_x}{y}\right)\ez = \frac{1}{2}\, \rot \Vec{v},
\end{equation}
où l'on a définit le champ vectoriel $\rot \Vec{v}$, appelé \imp{rotationnel}, pour un champ de vitesse $\Vec{v}$ bidimensionnel.
}

\begin{prop}[Interprétation du rotationnel]\label{prop:}
Le \imp{rotationnel} du champ de vitesse, noté $\Rot{\Vec{v}}$, caractérise la rotation locale des particules de fluide.
\end{prop}

% \begin{rappel}{Opérateur rotationnel}


% \end{rappel}



% \begin{exemple}{Écoulement et rotation solide}
% Considérons un écoulement plan stationnaire de champ de vitesse
% $$
% \Vec{v} = \omega r \et
% $$
% en coordonnées polaires $(r,\theta)$, avec $\omega$ une constante homogène à une vitesse angulaire. Alors le rotationnel vaut
% $$
% \Rot{\Vec{v}} = \omega \ez,
% $$
% qui est constant. Cela montre que le fluide tourne d’un seul bloc, comme un solide. La rotation se fait autour de l’axe $(Oz)$ dans le sens direct si $\omega>0$, et s’effectue à la vitesse angulaire $\omega$.
%
% \end{exemple}

 \begin{definition}[Écoulement irrotationnel]\label{def:irro}
 Un écoulement est dit \imp{irrotationnel} lorsque
$
 \Rot{\Vec{v}} = \Vec{0}.
 $
 \end{definition}
% \begin{definition}\label{def:irro}
% Un écoulement est dit \imp{irrotationnel}, ou potentiel, lorsque, de façon équivalente,
% \begin{equation*}
% \Rot{v} = \Vec{0} \text{~ou~} \Vec{v}=\Grad{\phi}.
% \end{equation*}
% \end{definition}

\begin{remarque}
  Par un calcul d'analyse vectorielle, on peut montrer que l'accélération convective se met sous la forme
% \begin{equation}
\eqbox[]{
  % \begin{split}
 (\Vec{v} \cdot \grad) \Vec{v} =  \Grad{\left(\frac{v^2}{2}\right)}+(\Rot{\Vec{v}})\pv\Vec{v}.
  % \end{split}
  }
% \end{equation}
Dans un écoulement irrotationnel, on a donc 
\begin{equation}
  (\Vec{v} \cdot \grad) \Vec{v} =  \Grad{\left(\frac{v^2}{2}\right)} \Rar \Dt{\Vec{v}} = \dr{\Vec{v}}{t} + \Grad{\left(\frac{v^2}{2}\right)}.
\end{equation}
\end{remarque}

%, pour lesquels $\Rot{\Vec{v}}=0$.

\parag{Potentiel des vitesses}

Les propriétés du rotationnel entraîne que $\Rot{\Grad{\phi}}=\Vec{0}$ pour tout champ scalaire $\phi(M,t)$.
La réciproque est également vrai : si $\Rot{\Vec{v}}(M,t) = \Vec{0}$ en tout point de l'espace, il existe un potentiel $\phi(M,t)$ tel que $\Vec{v} = \Grad{\phi}$.
\noindent Le champ scalaire $\phi(M,t)$ est appelé le \imp{potentiel des vitesses}.

\begin{prop}[Écoulement potentiel]\label{prop:}
  Un écoulement est \imp{irrotationnel} si et seulement s'il est \imp{potentiel}, i.e. 
  \begin{equation*}
    \Rot{\Vec{v}} = \Vec{0} \Leftrightarrow \exists~\phi(M,t), \Vec{v} = \Grad{\phi}.
  \end{equation*}
\end{prop}
\begin{remarque}
  Le choix du signe dans la définition de $\phi$ est conventionnel. Cette convention standard en mécanique des fluides  s'oppose aux conventions choisies en mécanique, où une force conservative s'écrit $\Vec{F}=-\Grad{E_{\rm{pot}}}$, ou encore en électrostatique, où le champ électrique vérifie $\Vec{E}=-\Grad{V}$. Attention, donc !
\end{remarque}

\begin{exemple}{Exemples d’écoulements potentiels}
  Associer chacun des écoulements bidimensionnels suivants à son potentiel des vitesses.
  \begin{fig}[label]
    \centering
    \begin{subfig}[label][0.45]
      \includegraphics[height=4cm]{flow_uniform.pdf}
    \end{subfig}
    \begin{subfig}[label][0.45]
      \includegraphics[height=4cm]{flow_source.pdf}
    \end{subfig}
    \begin{subfig}[label][0.45]
      \includegraphics[height=4cm]{flow_vortex.pdf}
    \end{subfig}
    \begin{subfig}[label][0.45]
      \includegraphics[height=4cm]{flow_dipole.pdf}
    \end{subfig}
    %\capt{}
  \end{fig}
  $$
    \begin{array}{rr}
      \phi = U x & (1) \\[5mm]
      \phi = \dfrac{Q}{2\pi} \ln r & (2) \\
    \end{array}
    \qquad
    \begin{array}{rr}
      \phi = -\dfrac{\Gamma \theta}{2\pi} & (3) \\[5mm]
      \phi = -\dfrac{K \cos \theta}{r} & (4)
    \end{array}
  $$
\end{exemple}
\examplesol{}


% \begin{remarque}
% Un cas particulier notable est celui d'un écoulement \imp{irrotationnel et incompressible}, auquel cas $\Dive{\Vec{v}}=\Dive{\grad{\phi}}=0$, et donc le potentiel vérifie l'équation de Laplace
% \begin{equation}
% \Delta\phi=0.
% \end{equation}
% \end{remarque}


%\begin{Code}
%ecoulement_sphere.py
%\end{Code}
%\code{ecoulement_sphere.py}
%\inputminted{python}{./Code/ecoulement_sphere.py}
%%\begin{figure}[h!]
%\centering
%\includegraphics[width=\textwidth]{code_lignes_courant.pdf}
%\caption{Code Python utilisé pour le tracé de la Fig.~\ref{fig:écoulement_sphère}.}
%\label{fig:écoulement_sphère_code}
%\end{figure}

%\subsection{L'écoulement uniforme}
%Dans ce cas, la vitesse $\Vec{v}$ est la même quel que soit le point du fluide considéré. Ainsi la dérivée convective de la vitesse $(\Vec{v}\cdot\Vec{\nabla})\Vec{v}$ est nulle et l'accélération particulaire se réduit au terme local $\frac{\partial \Vec{v}}{\partial t}$.

%
% \clearpage
% \begin{subappendices}
% \section{Déformations d'une particule de fluide en écoulement}
% \label{comp:déformation_carré}
% On cherche à décrire comment une particule de fluide est modifiée dans un écoulement quelconque. Pour simplifier, on considère l'écoulement bidimensionnel et on se place en coordonnées cartésiennes. On modélise alors une particule de fluide par une surface infinitésimale (un volume à 2D), que l'on peut prendre carrée pour simplifier, de côtés de longueurs $\d x$ et $\d y$ respectivement. Soit $M$ un point quelconque de la particule de fluide. On cherche à savoir comment $M$ est modifié au cours de l'écoulement. Calculons alors pour une petite déformation
% \begin{equation}
%   \begin{split}
%   \label{eq:comp_deformation_generale}
%   \Vec{A'M'}-\Vec{AM}&=\d\Vec{AM}=\left[\Vec{v}(M)-\Vec{v}(A)\right]\d t=\left[\frac{\partial \Vec{v}}{\partial x}\d x+\frac{\partial \Vec{v}}{\partial y}\d y\right]\d t\\
%   &=\left[\frac{\partial v_x}{\partial x}\d x\,\ex+\frac{\partial v_y}{\partial y}\d y\,\Vec{e_y}\right]\d t+\left[\frac{\partial v_x}{\partial y}\d y\,\ex+\frac{\partial v_y}{\partial x}\d x\,\Vec{e_y}\right]\d t.
%   \end{split}
%   \end{equation}
%   Nous allons étudier le rôle de chaque terme entre crochets séparément.
% \begin{figure}[h!]
% \centering
% \includegraphics[width=0.9\textwidth]{deformations_carre.pdf}
% \caption{Modification de la particule de fluide durant $\d t$ dans le cas (a) d'une déformation pure (écoulement divergent irrotationnel), puis dans le cas (b) d'une rotation pure (écoulement incompressible et tourbillonnaire).}
% \end{figure}
%
% \subsection{Compression et dilatation, divergence}
% Pour ce premier cas, on suppose que $\Vec{v}$ est de la forme
% \begin{equation}
% \Vec{v}=v_{x}(x)\,\ex+v_{y}(y)\,\Vec{e_y},
% \end{equation}
% de sorte que seul le premier crochet de l'Éq.~\eqref{eq:comp_deformation_generale} est non nul.
% Si le point M est en B, alors $\d y=0$ et donc $\d \Vec{AB}=\frac{\partial v_x}{\partial x}\d x\d t\,\ex$, qui est parallèle à $\Vec{AB}$. Autrement dit, l'allongement relatif horizontal de la particule de fluide vaut $\frac{\d AB}{AB}=\frac{\partial v_x}{\partial x}\d t$ puisque $AB=\d x$.
% Exactement de la même manière, si le point M est en D, alors $\d x=0$ et il vient que l'allongement relatif vertical de la particule de fluide vaut $\frac{\d AD}{AD}=\frac{\partial v_y}{\partial y}\d t$.
%
% Calculons alors la variation relative de surface (volume bidimensionnel) de la particule de fluide. La surface originale vaut $S=\mathrm{AB}\times\mathrm{AD}=\d x\,\d y$, et donc en utilisant la différentielle d'un produit
% \begin{equation}
% \begin{split}
% \frac{\d S}{S}&=\frac{\d(\mathrm{AB})\times \mathrm{AD}}{S}+\frac{\d(\mathrm{AD})\times \mathrm{AB}}{S}=\left(\frac{\partial v_x}{\partial x}+\frac{\partial v_y}{\partial y}\right)\d t
% \end{split}
% \end{equation}
% Ce résultat se généralise instantanément à trois dimensions et on obtient que
% \begin{equation}
% \frac{\d V}{V}=\Dive{\Vec{v}}\,\d t,
% \end{equation}
% qui est l'Éq.~\eqref{eq:interp_div_variation_volume} que nous avions déjà obtenu en Sec.~\ref{subsec:alternative_cons_masse} pour l'interprétation de la divergence.
% \subsection{Rotation sans déformation, rotationnel}
% Cette fois on suppose que $\Vec{v}$ est de la forme
% \begin{equation}
% \Vec{v}=v_{x}(y)\,\ex+v_{y}(x)\,\Vec{e_y},
% \end{equation}
% de sorte que seul le deuxième crochet de l'Éq.~\eqref{eq:comp_deformation_generale} est non nul.
% Si le point M est en B, $\d y=0$ conduit à $\d \Vec{AB}=\frac{\partial v_y}{\partial x}\d x\d t\,\Vec{e_y}$, qui est orthogonal à $\Vec{AB}$. Donc la partie horizontale de la particule de fluide tourne autour de $\Vec{e_z}$, d'un angle $\d\alpha$ petit que l'on peut exprimer comme
% \begin{equation}
% \tan(\d\alpha)=\frac{\frac{\partial v_y}{\partial x}\d x\d t}{\d x}\simeq \d\alpha
% \end{equation}
% et donc la vitesse angulaire de rotation autour de $\Vec{e_z}$ s'écrit
% \begin{equation}
% \frac{\d\alpha}{\d t}=\dot{\alpha}=\frac{\partial v_y}{\partial x}.
% \end{equation}
% Exactement de la même manière, si le point M est en D, alors $\d x=0$ et on montre qu'une rotation à la vitesse angulaire
% \begin{equation}
% \dot{\beta}=-\frac{\partial v_x}{\partial y}
% \end{equation}
% se met en place (le signe - provient du fait que $\d\Vec{AD}$ est selon $-\ex$, voir schéma).
%
% Que se passe-t-il si $\dot{\alpha}=\dot{\beta}$ ? Dans ce cas, il y a une rotation pure de la particule de fluide, à la vitesse angulaire
% \begin{equation}
% \begin{split}
% \dot{\Omega}&=\dot{\alpha}=\dot{\beta}=\frac{\dot{\alpha}+\dot{\beta}}{2}=\frac{1}{2}\left(\frac{\partial v_y}{\partial x}-\frac{\partial v_x}{\partial y}\right)\\
% &=\frac{1}{2}\left[\begin{pmatrix} \frac{\partial}{\partial x} \\[2mm] \frac{\partial}{\partial y} \\[2mm] / \end{pmatrix}\pv \begin{pmatrix} v_x \\[2mm] v_y \\[2mm] / \end{pmatrix}\right]\cdot\Vec{e_z}=\frac{1}{2}\Rot{\Vec{v}}\cdot\Vec{e_z}
% \end{split}
% \end{equation}
% où dans la dernière ligne nous illustrons la généralisation à trois dimension. La quantité $\Vec{\Omega}=\frac{1}{2}\Rot{\Vec{v}}$ représente donc la rotation locale sans déformation de la particule de fluide.
% % Il s'agit du vecteur tourbillon défini en Sec.~\ref{subsec:ecoulement_tourbillonnaire}.
%
% \end{subappendices}

\end{document}
